\section{Introduction}

Estimating the geographic origin of a photograph from its pixels alone has
progressed from retrieval over geotagged collections~\cite{hays2008im2gps},
through classification over discretised geographic
cells~\cite{weyand2016planet,pramanick2022translocator}, to contrastive
alignment between images and continuous GPS coordinates.
GeoCLIP~\cite{cepeda2023geoclip} exemplifies the latter: a
CLIP~\cite{radford2021clip} vision tower is aligned with a location encoder,
and inference reduces to nearest-neighbour retrieval over a gallery of
$10^{5}$ coordinates. Such models are accurate enough to be practically useful
and, for the same reason, consequential --- the capability supports both the
verification of news imagery and the inference of location from photographs
shared without geotags.

Despite this, the evidential basis of an individual prediction remains
largely opaque. It is widely assumed that these models exploit road
markings, vegetation, licence plates and shopfront text, but the supporting
evidence is predominantly anecdotal or derived from attention maps.
Whether raw attention weights explain predictions remains
contested~\cite{jain2019attention,wiegreffe2019attention}. Gradient-based
maps~\cite{selvaraju2017gradcam}, attention rollout~\cite{abnar2020rollout}
and Transformer relevance propagation~\cite{chefer2021transformer} assign
importance to regions, but do not directly test the output under an internal
intervention; saliency methods can also fail basic model- and data-randomisation
checks~\cite{adebayo2018sanity}. Input-space perturbation methods instead mask
pixels~\cite{petsiuk2018rise} or replace image regions~\cite{goyal2019counterfactual},
but may confound model sensitivity with the distribution shift caused by an
altered image. The practical question here is how a prediction changes when a
specific attention pathway is manipulated while pixels and weights stay fixed.
Such analysis is established in language models, where interventions identify
internal routes associated with model behaviour~\cite{vig2020causal,meng2022rome,geva2023dissecting}.
We transfer that approach to image
geolocation and make its geographic effect directly observable.

We present GeoProbe, whose contributions are threefold. First, an
interactive system applying region-conditioned, per-layer attention
intervention to the GeoCLIP vision tower at inference time, requiring
neither training, gradients nor architectural modification, and presenting
the unmodified and intervened predictions jointly on a map. Second, an
object-proposal workflow in which WeDetect-Uni~\cite{fu2026wedetect}
supplies category-agnostic foreground regions without requiring a prompt for
each image, enabling the same intervention to be evaluated at dataset scale.
Third, a quantitative study on Img2GPS3k~\cite{vo2017revisiting} that uses a
discovery/holdout protocol to test single layers and multi-layer
combinations. The study quantifies output sensitivity while showing that
apparently favourable search results need not generalise to unseen images.
